\section{Identities for planar sections}
\label{slice:section}

Coordinate monotonicity lets us describe the boundary through planar
sections. The identities below convert their perimeters and area changes
into bounds on the total boundary area at the two transition volumes.

Choose an axial direction \(z\) in the reflected cube \(Q\), and write
\[
 E_z=\{(x,y)\in[0,1]^2:(x,y,z)\in E\},\qquad
 v(z)=|E_z|,\quad s=v(0),\quad r=v(1).
\]
Let \(p(z)\) be the relative perimeter of \(E_z\) in the square.
In slice integrals, \(\partial E_z\) denotes the separating curve
\(\{(x,y)\in[0,1]^2:(x,y,z)\in\Sigma\}\), including its endpoints;
portions of the square boundary are not counted. The surface coarea
Jacobian for the height function is
\[
 \alpha=|\nabla_\Sigma z|=\sqrt{1-N_z^2}.
\]
Its mean square over the surface patch will be denoted by
\[
 b_z=\frac1A\int_{\Sigma\cap Q}(1-N_z^2)\,dS.
\]
All section identities at interior heights are initially understood
at regular levels.

\begin{lemma}[Regular levels and area coordinates]\label{slice:regularity}
There are only finitely many exceptional axial heights.
The function \(v\) is continuous, nonincreasing, absolutely continuous,
and smooth off those heights. The end levels are regular, with a regular
collar whenever nonempty. For every nonnegative measurable function
\(\gamma\),
\begin{equation}
 \int_0^1 \gamma(v(z))(-v'(z))\,dz
       =\int_r^s\gamma(w)\,dw.
 \label{slice:substitution}
\end{equation}
\end{lemma}

\begin{proof}
A critical point of the height on \(\Sigma\) has \(N=\pm e_z\).
By \eqref{geo:strict-normal}, both \(x,y\) then belong to \(\{0,1\}\).
Such a point lies on an axial edge, which is transverse to the surface.
These intersections are isolated and, by compactness and the exclusion
of vertices, finite.

The boundary faces of the square introduce no further exceptional
heights. On \(x=c\), \(c\in\{0,1\}\), one has \(N_x=0\);
a tangent to the face intersection is
\(e_x\times N=(0,-N_z,N_y)\).
Its height derivative can vanish only if \(N_y=0\), forcing \(y=0\)
or \(1\), and hence an already-counted axial edge. Interchanging
\(x,y\) treats the other faces. Equivalently, reflect in these two
coordinates and take sections in a transverse two-torus.

At the end faces \(N_z=0\), so the height restricted to the reflected
surface is regular there. Compactness supplies a regular collar;
an empty end section has an empty collar.
At any interior height the intersection with \(\Sigma\) is a union
of regular curve pieces and finitely many critical points, of zero
planar area. Any jump of the nested section areas would be supported
on that boundary intersection. Thus \(v\) is continuous. On intervals
of regular levels it is smooth by smooth transport of the compact
section curves.

Monotonicity makes \(Dv\) a finite measure. Its singular part can be
supported only on the finite set of exceptional heights. Continuity
rules out atoms, and therefore this singular part is zero.
This proves absolute continuity without first assuming integrability
of \(1/\alpha\).

For \eqref{slice:substitution}, push forward the measure
\((-v'(z))\,dz\) under \(v\). For \(r<a<b<s\), the inverse image of
\((a,b)\) is an interval, up to endpoint plateaus, and absolute
continuity gives it measure \(b-a\). A level set has zero measure for
this pushforward: on any plateau \(v'=0\) almost everywhere.
The pushforward is therefore Lebesgue measure on \((r,s)\).
The assertion follows first for interval indicators and then for all
nonnegative measurable \(\gamma\). No regularity of the inverse of
\(v\) is required.
\end{proof}

\begin{lemma}[Flux and coarea]\label{slice:identities}
With the preceding notation,
\begin{align}
 \int_{\partial E_z}\alpha\,d\ell
     &=b_z A+h(v(z)-V),\notag\\
 p(0)&=b_zA+h(s-V),&
 p(1)&=b_zA+h(r-V),
 \label{slice:flux}\\
 s-r&=\int_{\Sigma\cap Q}N_z\,dS
       \le A\sqrt{1-b_z},
 \label{slice:projection}\\
 -v'(z)&=\int_{\partial E_z}\frac{N_z}{\alpha}\,d\ell,&
 A&=\int_0^1\int_{\partial E_z}\frac1\alpha\,d\ell\,dz.
 \label{slice:coarea}
\end{align}
Moreover \(b_x+b_y+b_z=2\), and \(b_i<1\) in every direction.
\end{lemma}

\begin{proof}
Apply first variation to \(X=\eta(z)e_z\), where
\(\eta\in C_c^\infty(0,1)\). Orthogonality to the cube faces eliminates
the boundary term; the same computation may be made after reflection.
The outward/sum-curvature convention and the divergence theorem give
\[
 \int_{\Sigma\cap Q}\eta'(z)(1-N_z^2)\,dS
       =h\int_{\Sigma\cap Q}\eta(z)N_z\,dS
       =h\int_0^1\eta'(z)v(z)\,dz.
\]
Coarea converts the first integral to
\(\int\eta'(z)\int_{\partial E_z}\alpha\,d\ell\,dz\).
Consequently
\(\int_{\partial E_z}\alpha\,d\ell-hv(z)\) is constant.
Averaging over the unit height interval identifies it as \(b_zA-hV\).
Continuity on regular intervals gives the identity at every regular
level. At each end, \(\alpha=1\), and the regular collar gives the
two endpoint formulas in \eqref{slice:flux}.

The divergence theorem for \(e_z\) gives
\(\int_{\Sigma\cap Q}N_z\,dS=s-r\). Since
\(\int N_z^2\,dS=A(1-b_z)\), Cauchy--Schwarz yields
\eqref{slice:projection}.

At a regular height, the outward planar normal velocity of the slice
curve is \(-N_z/\alpha\). The area-variation formula gives the first
identity in \eqref{slice:coarea}. Coarea on the surface gives the second;
the finite critical-point set has surface area zero, so no term is
lost there. These integrals remain valid if their integrands become
unbounded as an exceptional height is approached.

Finally \(\sum_i(1-N_i^2)=2\) pointwise, giving the sum of the \(b_i\).
The strict sign of \(N_i\) on a nonempty interior patch gives
\(\int N_i^2>0\), hence \(b_i<1\).
\end{proof}
